\documentclass[12pt]{exam}

\usepackage[utf8]{inputenc}  
\usepackage[portuguese]{babel}   

\usepackage[ruled, vlined, algonl, portuguese]{algorithm2e}

\newcommand{\linesnumbered}{\LinesNumbered}
\newcommand{\dontprintsemicolon}{\DontPrintSemicolon}
\newcommand{\incmargin}{\IncMargin}
\newcommand{\decmargin}{\DecMargin}

\usepackage{mathtools}

\usepackage{amssymb}
\usepackage{amsmath}
\usepackage[shortlabels]{enumitem}

\def\le{\leqslant}\def\leq{\le}
\def\ge{\geqslant}\def\geq{\ge}
\newcommand{\floor}[1]{\left\lfloor{#1}\right\rfloor}
\newcommand{\ceil}[1]{\left\lceil{#1}\right\rceil}
\newcommand{\pare}[1]{\left({#1}\right)}
\newcommand{\set}[1]{\left\{{#1}\right\}}
\newcommand{\range}[2]{\set{{#1},\dots,{#2}}}
\newcommand{\dist}{\mathrm{dist}}
\newcommand{\cluster}{\mathrm{cluster}}
\newcommand{\merge}{\mathrm{merge}}

\pagestyle{empty}

\title{Lista 02 de Álgebra Linear}
\date{}



\begin{document}


\maketitle


\vspace{3em}


\begin{questions}

    \question Normalize os seguintes vetores:
    \begin{parts}
        \part (0, 1)
        \part (1, 0)
        \part (1, 1)
        \part (3, 4)
        \part (6, 8)
    \end{parts}


    \question Considere os vetores (1, 2) e (0, 1).
    Prove que os seguintes vetores são 
    combinação linear destes vetores.

    \begin{parts}
        \part (0, 0)
        \part (0, 1)
        \part (1, 0)
        \part (1, 1)
        \part (2, 2)
    \end{parts}

    \question Considere os vetores (1, 2) e (2, 4).
    Verifique \textbf{SE} os seguintes vetores são
    combinação linear destes vetores.

    \begin{parts}
        \part (0, 0)
        \part (0, 1)
        \part (1, 0)
        \part (1, 1)
        \part (2, 2)
    \end{parts}

\break

\question Considere os vetores 
    (1, 1, 1), 
    (2, 3, 2) e
    (3, 2, 4).
Prove que os seguintes vetores são 
combinação linear destes vetores.

    \begin{parts}
        \part (2, 1, 2)
        \part (2, 1, 1)
        \part (-1, 0, -1)
    \end{parts}



\question Usando o cálculo do produto escalar, 
calcule o cosseno do ângulo entre o vetor 
$(3, 4)$ e $(6, 8)$


\question Usando o cálculo do produto escalar, 
calcule o cosseno do ângulo entre o vetor 
$(2, 1)$ e $(-2, 4)$


\question Usando o cálculo do produto escalar, 
calcule o cosseno do ângulo entre o vetor 
$(3, 3)$ e $(0, 2)$


\question Descubra qual os valores de 
$\cos(45^\circ)$,  
$\cos(0^\circ)$ e  
$\cos(90^\circ)$  
usando o cálculo do produto escalar e 
os vetores $(1, 0)$, $(0, 1)$ e $(1, 1)$.

\end{questions}

\end{document}
